% Para preencher 

\newcommand{\ESCOLA}{<<Nome da Escola>>}
%\newcommand{\ESCOLA}{Escola Superior de Tecnologia e Gestão}
%\newcommand{\ESCOLA}{Escola Superior de Educação}
%\newcommand{\ESCOLA}{Escola Superior Agrária}
%\newcommand{\ESCOLA}{Escola Superior de Saúde}


%\newcommand{\CURSO}{Licenciatura/Mestrado/... em ...}

\newcommand{\TITULO}{<<Título>>}

% Se pretender subtítulo
\newcommand{\SUBTITULO}{Opcionalmente, colocar aqui sub-título com máximo de vinte palavras}
% Se NÃO pretender subtítulo retire o % (comentário) da linha seguinte e coloque na linha anterior
%\newcommand{\SUBTITULO}{} 

\newcommand{\NOMEALUNO}{<<Nome Completo Estudante>>}

\newcommand{\DECLARACAO}{<<Dissertação apresentada ao instituto politécnico de beja para obtenção do grau de mestre em xx, ramo de xx>>}

\newcommand{\LOCAL}{<<LOCAL>>}
\newcommand{\MES}{<<MES>>}
\newcommand{\ANO}{<<ANO>>}


\newcommand{\orientacao}{ORIENTAÇÃO}
\newcommand{\nomesorientacao}{Colocar o título e nome de ou dois orientadores, por exemplo "Professora Doutora Nome Completo ou Abreviado\\Eng. Nome Completo ou Abreviado".}

% se não existir supervisor (docente do IPBeja, com orientador noutra entidade, colocar comentários nas linhas seguintes
\newcommand{\supervisao}{SUPERVISÃO}
\newcommand{\nomesupervisao}{Colocar o título e nome do supervidor, por exemplo "Professora Doutora Nome Completo ou Abreviado"}
